

\section{Egri chiziqli harakat}


\begin{enumerate}

\item Disk 1 minutda 1200 marta aylanadi. Uning aylanish chastotasini toping (1/s).
\begin{enumerate}[label=\Alph*)]
\item 20
\item 1200
\item 60
\item 30
\end{enumerate}

\item Aylanish chastotasi 30 Hz bo‘lgan disk 5 s da necha marta aylanadi?
\begin{enumerate}[label=\Alph*)]
\item 20
\item 150
\item 6
\item 30
\end{enumerate}

\item Aylanish chastotasi 20 s$^{-1}$ bo‘lgan disk 1000 marta aylanishi uchun qancha vaqt kerak bo‘ladi?
\begin{enumerate}[label=\Alph*)]
\item 20
\item 0,02
\item 50
\item 1000
\end{enumerate}

\item Mushuk aylanani 20 s da 4 marta aylandi. Mushukning aylanish davrini toping (s).
\begin{enumerate}[label=\Alph*)]
\item 20
\item 80
\item 50
\item 5
\end{enumerate}

\item Sportchi yugurish maydonchasini 10 marta aylanishi uchun qancha vaqt kerak bo‘ladi (min)? Uning aylanish davri 42 s.
\begin{enumerate}[label=\Alph*)]
\item 10
\item 7
\item 42
\item 420
\end{enumerate}

\item Aylanish davri 0,785 s bo‘lgan jismning burchak tezligini toping (rad/s).
\begin{enumerate}[label=\Alph*)]
\item 1,57
\item 8
\item 3,14
\item 6,28
\end{enumerate}

\item Aylanish chastotasi 0,0796 s$^{-1}$ bo‘lgan maxovikning maksimal burchak tezligini toping (rad/s).
\begin{enumerate}[label=\Alph*)]
\item 10
\item 0,5
\item 4
\item 20
\end{enumerate}

\item Aylanish davri 3 s bo‘lgan jism 4 s da necha gradusga buriladi?
\begin{enumerate}[label=\Alph*)]
\item 270$^\circ$
\item 480$^\circ$
\item 300$^\circ$
\item 360$^\circ$
\end{enumerate}

\item Quyidagilarni qaysi biri $12\cdot10^3$ sekund vaqt miqdorini ifodalaydi?
\begin{enumerate}[label=\Alph*)]
\item 2 soat–40 min
\item 3 soat–20 min
\item 4 soat–15 min
\item 4 soat–40 min
\end{enumerate}

\item Yukni 0,4 m/s tezlik bilan ko‘tarishda diametri 16 sm bo‘lgan chig‘ir barabanining aylanish chastotasi qanday bo‘lishining toping (s$^{-1}$).
\begin{enumerate}[label=\Alph*)]
\item 0,8
\item 16
\item 0,4
\item 0,2
\end{enumerate}

\item Radiusi 0,25 m bo‘lgan g‘ildirak 12 m/s tezlik bilan g‘ildiramoqda. G‘ildirakning burchak tezligini toping (rad/s).
\begin{enumerate}[label=\Alph*)]
\item 12
\item 24
\item 36
\item $12\pi$
\end{enumerate}

\item Sharcha radiusi 1 m bo‘lgan aylana bo‘ylab doimiy 4 m/s tezlik bilan harakatlanmoqda. Uning aylanish davri qanday (s)?
\begin{enumerate}[label=\Alph*)]
\item 1,57
\item 3,14
\item 4,71
\item 6,28
\end{enumerate}

\item Avtomobilning tezligi 72 km/soat, g‘ildiragining diametri 62 sm. G‘ildirakning burchak tezligini toping (rad/s larda).
\begin{enumerate}[label=\Alph*)]
\item 64,5
\item 62,5
\item 62
\item 32
\end{enumerate}

\item Elektovoz g‘ildiragining diametri 1 m bo‘lib, u 1 minutda 300 marta aylansa, poezdning tezligi necha m/s bo‘ladi? ($\pi=3$ deb hisoblang).
\begin{enumerate}[label=\Alph*)]
\item 10
\item 12
\item 15
\item 20
\end{enumerate}

\item Aylana bo‘ylab tekis harakat qilayotgan moddiy nuqtaning chiziqli tezligi 2 m/s va burchak tezligi 5 rad/s bo‘lsa, markazga intilma tezlanishi qanday (m/s$^2$)?
\begin{enumerate}[label=\Alph*)]
\item 0,4
\item 2
\item 2,5
\item 5
\end{enumerate}

\item Poezd egrilik radiusi 800 m bo‘lgan burilishda 20 m/s tezlik bilan harakatlanayotganining markazga intilma tezlanishi qanday (m/s$^2$)?
\begin{enumerate}[label=\Alph*)]
\item 10
\item 1
\item 0,2
\item 0,5
\end{enumerate}

\item Velosipedchi egrilik radiusi 48 m bo‘lgan burilishda 12 m/s tezlikda harakatlanmoqda. Markazga intilma tezlanishini toping (m/s$^2$).
\begin{enumerate}[label=\Alph*)]
\item 4
\item 3
\item 0,25
\item 0,5
\end{enumerate}

\item Lokomotiv yo‘lining radiusi 750 m bo‘lgan burilish joyidan 54 km/soat tezlik bilan o‘tmoqda. Tezligi 2 marta kamaysa, lokomotivning markazga intilma tezlanishi qanday o‘zgaradi?
\begin{enumerate}[label=\Alph*)]
\item 3 marta kamayadi
\item 4 marta kamayadi
\item 2 marta kamayadi
\item 4 marta ortadi
\end{enumerate}

\item Chekka nuqtasining markazga intilma tezlanishi 49,3 m/s$^2$ bo‘lgan diskning radiusi 5 sm. Diskning aylanish chastotasini toping (Hz).
\begin{enumerate}[label=\Alph*)]
\item 5
\item 10
\item 3
\item 6
\end{enumerate}

\item Radiusi 10 sm bo‘lgan aylana bo‘ylab tekis harakat qilayotgan moddiy nuqtaning aylanish davri 2 s bo‘lsa, uning markazga intilma tezlanishini toping (m/s$^2$).
\begin{enumerate}[label=\Alph*)]
\item 1,57
\item 0,57
\item 1,1
\item 1,05
\end{enumerate}

\item Aylana bo‘ylab tekis harakat qilayotgan jismning burchak tezligi 5,5 marta ortsa, chiziqli tezligi esa shuncha marta kamaymasa, jismning markazga intilma tezlanishi qanday o‘zgaradi?
\begin{enumerate}[label=\Alph*)]
\item 8,5 marta ortadi
\item 8,5 marta kamayadi
\item o‘zgarmaydi
\item 17 marta kamayadi
\end{enumerate}

\item Avtomobil 72 km/soat tezlik bilan harakatlanganda g‘ildiraklarining aylanish chastotasi 8 s$^{-1}$ bo‘lsa, avtomobil g‘ildiraklarining yo‘lga tegadigan nuqtalarining markazga intilma tezlanishini topilsin (m/s$^2$).
\begin{enumerate}[label=\Alph*)]
\item 10
\item 2
\item 15
\item 1
\end{enumerate}

\item Aylanish davri 0,2 s va radiusi 5 sm bo‘lgan disk chekka nuqtasining markazga intilma tezlanishini toping (m/s$^2$).
\begin{enumerate}[label=\Alph*)]
\item 9,9
\item 12,3
\item 15,7
\item 49,3
\end{enumerate}

\item Velosiped 10 m/s o‘zgarmas tezlik bilan harakatlanmoqda. Uning g‘ildiragi diametri 80 sm. G‘ildirakning pastki nuqtasi yerga nisbatan necha m/s$^2$ tezlanish bilan harakatlanadi?
\begin{enumerate}[label=\Alph*)]
\item 0
\item 1,25
\item 250
\item 125
\end{enumerate}

\item Jism aylana bo‘ylab tekis harakat qilmoqda. Agar aylana radiusi o‘zgarmasdan, jism tezligi 3 marta kamaysa, markazga intilma tezlanish qanday o‘zgaradi?
\begin{enumerate}[label=\Alph*)]
\item 9 marta kamayadi
\item 3 marta ortadi
\item 3 marta kamayadi
\item 9 marta ortadi
\end{enumerate}

\item $R_1$ va $R_2=4R_1$ radiusli aylana bo‘ylab harakatlanayotgan moddiy nuqtalarning aylanish davrlari teng bo‘lsa, ularning markazga intilma tezlanishlari uchun quyidagi munosabatlarning qaysi biri o‘rinli bo‘ladi?
\begin{enumerate}[label=\Alph*)]
\item $a_1=a_2$
\item $a_1=2a_2$
\item $a_2=4a_1$
\item $a_2=2a_1$
\end{enumerate}

\item Gidroturbina ishchi g‘ildiragining radiusi bug‘ turbinasinikiga qaraganda 8 marta katta, aylanish chastotasi esa 40 marta kichik. Turbina g‘ildiraklarining chetidagi nuqtalar tezlanishlarini taqqoslang.
\begin{enumerate}[label=\Alph*)]
\item 1:2
\item 1:20
\item 1:200
\item 1:400
\end{enumerate}

\item Moddiy nuqta aylana bo‘ylab 25 m/s tezlik bilan tekis aylanmoqda. Davrning to‘rtdan bir qismida moddiy nuqta tezligi o‘zgarishining moduli qanday (m/s)?
\begin{enumerate}[label=\Alph*)]
\item 0
\item 5
\item 25
\item $25\sqrt{2}$
\end{enumerate}

\item Moddiy nuqta aylana bo‘ylab 6 m/s tezlik bilan tekis aylanmoqda. Davrning to‘rtdan bir qismida moddiy nuqta tezligi o‘zgarishining moduli qanday (m/s)?
\begin{enumerate}[label=\Alph*)]
\item 0
\item 3
\item 12
\item $6\sqrt{2}$
\end{enumerate}

\item Agar aylanayotgan disk ustidagi radiuslar 1 sm ga farq qiladigan ikki nuqtaning chiziqli tezliklari orasidagi farq 0,314 m/s ga teng bo‘lsa, diskning aylanish chastotasini toping (Hz).
\begin{enumerate}[label=\Alph*)]
\item 6,28
\item 5
\item 3,14
\item 1
\end{enumerate}

\end{enumerate}



\begin{center}
\renewcommand{\arraystretch}{1.3}
\begin{tabular}{|c|c|c|c|c|c|c|c|c|c|}
\hline
1A & 2B & 3C & 4D & 5B & 6B & 7C & 8A & 9C & 10C \\
\hline
11A & 12B & 13B & 14B & 15D & 16A & 17B & 18D & 19C & 20A \\
\hline
21D & 22B & 23D & 24C & 25D & 26C & 27C & 28D & 29D & 30C \\
\hline
\end{tabular}
\end{center}













